\section{Introduction}
\label{sec:intro}

\subsection{The Case for Combination Search in Alzheimer’s Disease}

Alzheimer's disease (AD) is the most common cause of dementia. Its two hallmarks are plaques of amyloid-$\beta$ outside nerve cells and tangles of tau protein inside them, and in lesser cases: inflammation, synaptic loss, vascular damage and altered brain metabolism. Drugs aimed at a single target are shown to be ineffective against AD, and the approved anti-amyloid antibodies only modestly slow decline, so previous work suggests that several processes may need to be treated at once. The annual review of the AD drug-development pipeline reports that 35\% of the drugs now in AD clinical trials are \emph{repurposed}, meaning they are already approved for another indication, and that these drugs make up almost half of Phase~2 agents \cite{cummings2026pipeline}. Repurposed drugs give us a pool of approved candidates with known safety profiles, which makes a search over their combinations practical. Clinical trials of combinations are slow and expensive, because the number of trial arms grows with every drug added, so a computational method that prioritizes a short list of promising combinations can notably aid in generating hypotheses prior to these programs.

\subsection{A Network Rule for Drug Combinations}

We based our proposal under the framework of network medicine treating the cell's proteins as a graph, the human \emph{interactome}. Proteins linked to the same disease cluster together in this graph, forming a \emph{disease module} \cite{barabasi2011network,menche2015uncovering}. A drug acts on a few proteins, its \emph{targets}. Drugs whose targets are unusually close to the disease module are more likely to treat that disease \cite{guney2016network}.

For combinations, Cheng, Kov\'acs and Barab\'asi \cite{cheng2019combinations} looked at how the targets of two drugs relate to each other and to the disease module. Of the six possible patterns, only one went with combinations that actually work in patients: both drugs are close to the disease module, but their targets are in different parts of the network. They dub this observation as \emph{complementary exposure}.

\begin{comment}
\begin{plainwords}
Picture the interactome as a city map, with AD living in one neighborhood. Each drug is a delivery truck that stops at a few buildings. A good \emph{single} drug stops near the AD neighborhood. A good \emph{group} of drugs stops near the AD neighborhood, but each truck covers a different street, so they do not waste effort on the same buildings. Choosing the best team of, say, three trucks out of thirty is the problem this project attempts to solve.
\end{plainwords}
\end{comment}

\subsection{Motivation as to how this is a Computational Problem}

Choosing $k$ drugs out of $n$ means searching $\binom{n}{k}$ combinations. This is trivial for $n=20$, $k=3$ (1{,}140 combinations), but it grows quickly: $\binom{100}{5}\approx 7.5\times 10^{7}$, $\binom{1000}{5}\approx 8\times 10^{12}$, and $\binom{1000}{10}\approx 2.6\times 10^{23}$. We observe that the search consists of many independent, regular operations, and each of our three methods stresses a different part of the machine:
\begin{itemize}
\item \textbf{Exact search} scores every combination, and each score is independent of the others, so one GPU thread can check a slice of them. We expect it to be limited by decoding combination indices and reading the score table.
\item \textbf{Simulated annealing} runs many independent search chains. Every accepted step inside a chain adds one row of the score table to a length-$n$ array and subtracts another, which is what a \texttt{ymm} register (eight 32-bit integers) and \texttt{vpaddd}/\texttt{vpsubd} are for. Only accepted steps do this work, so we expect the speedup to depend on the acceptance rate, which tends to fall as the temperature drops.
\item \textbf{Quantum simulation} stores $2^n$ complex amplitudes and, for every gate, applies the same small matrix to every group of two or four of them. Each gate streams the whole state through memory, so we expect it to be limited by memory bandwidth, and memory capacity caps $n$ at about 30 qubits on our single GPU.
\end{itemize}
SIMD instructions (AVX2) and GPUs (CUDA) suit this kind of regular, data-parallel work, but not every part speeds up equally. Part of our contribution is measuring where the speedup comes from and where it stops, and comparing all three methods in the range of $n$ where all of them can run.

\subsection{The Quantum-based Addition}
Extending beyond traditional computing. our scoring function has single terms and pair terms only, which is the shape of a Quadratic Unconstrained Binary Optimization (QUBO) problem \cite{glover2019tutorial,lucas2014ising}. QUBOs are the standard input for quantum optimization algorithms such as QAOA \cite{farhi2014qaoa}. We have no quantum computer, but NVIDIA cuQuantum \cite{bayraktar2023cuquantum} simulates one with an NVIDIA GPU that we are able to access. Because our instances are small enough to solve exactly, we can measure precisely how close simulated QAOA gets to the true best answer. A recent study suggests network-based drug-combination prediction as a QUBO and solved it with simulated quantum annealing for four other diseases \cite{ramos2025quantum}; we found no study that applies QAOA to AD drug combinations.

\subsection{Expected Deliverables}

\begin{enumerate}
\item A documented and checked scoring function (QUBO) for AD drug combinations, built from publicly-accessible drug data.
\item An exact GPU search over all valid combinations, which gives the reference answer.
\item One simulated-annealing algorithm written in three approaches (sequential scalar C++, AVX2 assembly, CUDA), with identical results, and an execution-time comparison.
\item QAOA, simulated with CUDA-Q and cuStateVec, measured against the exact answer.
\end{enumerate}.

\subsection{Relationship to Sister Work: Network-Screening project}

This project replaces the optional Module 5 of the sister project and can run independently. It needs two things for each candidate drug: its closeness score $z$ and its list of targets. It can take these from the sister project's output, or compute them for a hand-picked candidate list with a small script of its own (Decision D1, \Cref{sec:decisions}). The two projects agree on one table format (\Cref{app:data}), so either source can be plugged in without changing the underlying solvers.